\documentclass[11pt]{article}
\usepackage[T1]{fontenc}
\usepackage[utf8]{inputenc}
\usepackage{lmodern}
\usepackage[margin=1in]{geometry}
\usepackage{microtype}
\usepackage{amsmath,amssymb,amsthm,mathtools}
\usepackage{enumitem}
\usepackage{xcolor}
\usepackage{hyperref}
\usepackage{booktabs}
\usepackage{array}
\hypersetup{colorlinks=true,linkcolor=blue!45!black,citecolor=blue!45!black,urlcolor=blue!45!black}

\setlength{\parindent}{0pt}
\setlength{\parskip}{0.55em}
\setlist[itemize]{topsep=0.3em,itemsep=0.15em,leftmargin=1.5em}
\setlist[enumerate]{topsep=0.3em,itemsep=0.15em,leftmargin=1.8em}

\newtheorem{theorem}{Theorem}[section]
\newtheorem{lemma}[theorem]{Lemma}
\newtheorem{proposition}[theorem]{Proposition}
\newtheorem{corollary}[theorem]{Corollary}
\theoremstyle{definition}
\newtheorem{definition}[theorem]{Definition}
\newtheorem{remark}[theorem]{Remark}
\newtheorem{example}[theorem]{Example}

\newcommand{\ALCIRCC}{\mathsf{ALCI\_RCC5}}
\newcommand{\cl}{\operatorname{cl}}
\newcommand{\tp}{\operatorname{tp}}
\newcommand{\Roles}{\mathsf{R}}
\newcommand{\Eq}{\mathsf{EQ}}
\newcommand{\Dr}{\mathsf{DR}}
\newcommand{\Po}{\mathsf{PO}}
\newcommand{\Pp}{\mathsf{PP}}
\newcommand{\Ppi}{\mathsf{PPI}}
\newcommand{\inv}{\operatorname{inv}}
\newcommand{\Succ}{\operatorname{Succ}}
\newcommand{\repr}{\operatorname{repr}}
\newcommand{\occ}{\operatorname{Occ}}
\newcommand{\comp}{\operatorname{comp}}

\title{\textbf{On the Failure of Finite-Width Extraction for $\ALCIRCC$}\\[0.4em]
\large A counterexample to the $\mathsf{FW}(C,N)$ conjecture\\
and its implications for the contextual tableau approach}
\author{Claude (Anthropic)\\
\small Prompted by Michael Wessel}
\date{March 2026}

\begin{document}
\maketitle

\begin{center}
\fbox{\parbox{0.92\textwidth}{%
\textbf{Disclaimer.}\;
This document was generated by Claude (Anthropic), prompted by Michael
Wessel.  The results have not been independently verified.  We invite
scrutiny and corrections.}}
\end{center}

\begin{abstract}
A recent contextual tableau calculus for $\ALCIRCC$
\cite{GPT2026} reduces the decidability question to the
\emph{finite-width extraction} conjecture $\mathsf{FW}(C,N)$: every
satisfiable concept~$C$ admits, at some computable width bound~$N(C)$,
a finite closed family of contextual local states.  We show that
$\mathsf{FW}(C,N)$ is \textbf{false} for every~$N$ when $C$ is the
concept $(\exists\Pp.\top)\sqcap(\forall\Pp.\exists\Pp.\top)$, which
is satisfiable only in infinite models.  The argument is short: the
$\Pp$-relation in any local state is a finite strict partial order,
hence has maximal elements, whose existential $\Pp$-demands cannot be
met by the recentering mechanism.  The counterexample shows that the
contextual tableau framework, as formulated, cannot serve as a
complete decision procedure for $\ALCIRCC$.  The decidability of
$\ALCIRCC$ remains open.
\end{abstract}

%% ============================================================
\section{Introduction}

The description logic $\ALCIRCC$ extends $\mathcal{ALCI}$ with a role
box derived from the RCC5 composition table.  Wessel~\cite{Wessel2003}
introduced the logic and left decidability open, noting that
$\ALCIRCC$ lacks both the finite model property and the tree model
property.  In particular, the concept
\begin{equation}\label{eq:inf}
C_\infty \;=\; (\exists\Pp.\top)\sqcap(\forall\Pp.\exists\Pp.\top)
\end{equation}
is satisfiable only in infinite models \cite{Wessel2003,Claude2026}.

A 2026 discussion draft~\cite{Claude2026} proposed a quasimodel
approach but identified an \emph{extension gap} in the completeness
direction.  A companion note~\cite{GPT2026} proposed a
\emph{contextual tableau calculus} that avoids the extension gap by
re-designing the shape of a tableau node: each node is a finite
contextual local state---a bounded-width atomic RCC5 network of
remembered items with witness assignments and recentering maps.  The
note proves full \emph{soundness} (every open tableau graph unfolds
into a genuine model) but isolates completeness as an open conjecture:

\begin{quote}
$\mathsf{FW}(C,N)$: \emph{Every model of $C$ contains, after
quotienting by a finite descriptor equivalence, a closed family of
local states of width at most~$N$.}
\end{quote}

If $\mathsf{FW}(C,N(C))$ holds for some computable bound $N(C)$, then
combining it with soundness yields decidability \cite[Proposition~7.1]{GPT2026}.

\medskip
The purpose of this note is to show that $\mathsf{FW}(C_\infty,N)$
fails for \textbf{every}~$N$.  The argument is elementary and rests on
a single observation: $\Pp$ is a strict partial order, and finite
strict partial orders have maximal elements.

%% ============================================================
\section{Recollection of the framework}\label{sec:recap}

We recall only what is needed for the counterexample; see
\cite{GPT2026} for full definitions.

\subsection{Local states and closed families}

Fix a concept $C$ in negation normal form and let
$\cl(C)$ denote its closure under subformulas.

\begin{definition}[local state, \protect{\cite[Definition~3.1]{GPT2026}}]
A \emph{local state} of width at most~$N$ is a tuple
$M=(I_M,0_M,\tp_M,\rho_M,w_M)$ where:
\begin{enumerate}[label=(\alph*)]
\item $I_M$ is a finite nonempty set of \emph{items} with
  $|I_M|\le N$ and a distinguished \emph{center} $0_M\in I_M$.
\item $\tp_M:I_M\to 2^{\cl(C)}$ assigns Hintikka types.
\item $\rho_M:I_M\times I_M\to\Roles$ is a total, path-consistent
  atomic RCC5 labeling with $\rho_M(i,i)=\Eq$ and
  $\rho_M(j,i)=\inv(\rho_M(i,j))$.
\item For every $\exists R.D\in\tp_M(0_M)$, a designated witness
  $w_M(\exists R.D)\in I_M$ with
  $\rho_M(0_M,w_M(\exists R.D))=R$ and $D\in\tp_M(w_M(\exists R.D))$.
\end{enumerate}
The state satisfies universal propagation~(L1), transitive
propagation~(L2), and $\Eq$-synchrony~(L3).
\end{definition}

\begin{definition}[closed family / open tableau graph, \protect{\cite[Definitions~4.1, 5.1]{GPT2026}}]
A \emph{closed family of width~$N$} for $C$ consists of:
\begin{itemize}
\item a finite set $\mathcal{F}$ of saturated local states of width
  $\le N$, with a distinguished root state containing $C$ in its
  center type;
\item for each $M\in\mathcal{F}$ and each $i\in I_M$, a successor
  state $\Succ(M,i)\in\mathcal{F}$;
\item for each such edge, a total \emph{recentering map}
  $u_{M,i}:I_{\Succ(M,i)}\to I_M$ satisfying the axioms:
  \begin{enumerate}[label=(R\arabic*)]
  \item $u_{M,i}(0_{\Succ(M,i)})=i$;
  \item $\tp_{\Succ(M,i)}(a)=\tp_M(u_{M,i}(a))$ for all $a$;
  \item $\rho_{\Succ(M,i)}(a,b)=\rho_M(u_{M,i}(a),u_{M,i}(b))$ for all $a,b$;
  \item $\Succ(\Succ(M,i),a)=\Succ(M,u_{M,i}(a))$ for all $a$.
  \end{enumerate}
\end{itemize}
\end{definition}

The key result of \cite{GPT2026} is:

\begin{theorem}[Soundness, \protect{\cite[Corollary~6.6]{GPT2026}}]
If a closed family of width~$N$ for $C$ exists, then $C$ is
satisfiable in $\ALCIRCC$ under strong $\Eq$ semantics.\qed
\end{theorem}

The conditional completeness says: if $\mathsf{FW}(C,N)$ holds, then
$C$ is satisfiable \emph{if and only if} a closed family of width~$N$
exists.

%% ============================================================
\section{The counterexample}\label{sec:counter}

\subsection{The concept and its models}

Consider
\[
C_\infty = (\exists\Pp.\top)\sqcap(\forall\Pp.\exists\Pp.\top).
\]
As observed by Wessel~\cite{Wessel2003} and
Claude~\cite[Remark~2.5]{Claude2026}, $C_\infty$ is satisfiable only
in infinite models.  A canonical model is the $\Pp$-chain
$\mathcal{I}$ with domain $\Delta^{\mathcal{I}}=\{d_0,d_1,d_2,\dots\}$ and
\[
\rho(d_i,d_j) =
\begin{cases}
\Pp & \text{if } i<j,\\
\Eq & \text{if } i=j,\\
\Ppi & \text{if } i>j.
\end{cases}
\]
Every $d_n$ satisfies $\exists\Pp.\top$ (witnessed by $d_{n+1}$) and
$\forall\Pp.\exists\Pp.\top$ (every $\Pp$-successor $d_m$ with $m>n$
has its own $\Pp$-successor $d_{m+1}$).

\subsection{The type forced by transitive propagation}

Let $\tau$ be the Hintikka type that contains
$\exists\Pp.\top$ and $\forall\Pp.\exists\Pp.\top$.

\begin{lemma}\label{lem:type-propagation}
In any closed family for $C_\infty$, every item $i$ of every state $M$
that satisfies $\rho_M(0_M,i)=\Pp$ must have
$\exists\Pp.\top\in\tp_M(i)$ and
$\forall\Pp.\exists\Pp.\top\in\tp_M(i)$.
\end{lemma}

\begin{proof}
The center $0_M$ has $\forall\Pp.\exists\Pp.\top\in\tp_M(0_M)$ (since
$C_\infty\in\tp_M(0_M)$ for the root state, and for successor states
the type is inherited via~(R2)).

By universal propagation~(L1): since
$\forall\Pp.\exists\Pp.\top\in\tp_M(0_M)$ and $\rho_M(0_M,i)=\Pp$,
we get $\exists\Pp.\top\in\tp_M(i)$.

By transitive propagation~(L2): since $\Pp$ is transitive and
$\forall\Pp.\exists\Pp.\top\in\tp_M(0_M)$ with $\rho_M(0_M,i)=\Pp$,
we get $\forall\Pp.\exists\Pp.\top\in\tp_M(i)$.
\end{proof}

In the successor state $\Succ(M,i)$, item $i$ becomes the center.
By~(R2), the center of $\Succ(M,i)$ has the same type as item $i$ in
$M$.  If that type contains $\exists\Pp.\top$ and
$\forall\Pp.\exists\Pp.\top$, then $\Succ(M,i)$ is again a state
whose center has these formulas, and the lemma applies to its items as
well.

Hence: in any closed family for $C_\infty$, every state reachable from
the root has a center containing $\exists\Pp.\top$ and
$\forall\Pp.\exists\Pp.\top$.  Every $\Pp$-witness in such a state
also has these formulas in its type.

\subsection{The strict partial order obstruction}

\begin{lemma}\label{lem:spo}
In any local state $M$, the relation $\Pp$ on $I_M$ is a finite
strict partial order (irreflexive, asymmetric, transitive).
\end{lemma}

\begin{proof}
Irreflexivity: $\rho_M(i,i)=\Eq\ne\Pp$ for all $i$.

Asymmetry: if $\rho_M(i,j)=\Pp$ then $\rho_M(j,i)=\inv(\Pp)=\Ppi\ne\Pp$.

Transitivity: if $\rho_M(i,j)=\Pp$ and $\rho_M(j,k)=\Pp$, then by
the RCC5 composition table, $\rho_M(i,k)\in\comp(\Pp,\Pp)=\{\Pp\}$,
so $\rho_M(i,k)=\Pp$.
\end{proof}

\begin{lemma}\label{lem:maximal}
Every finite strict partial order has at least one maximal element.
\end{lemma}

\begin{proof}
Standard: if every element had a strict successor, iterating from any
element would produce an infinite chain, contradicting finiteness.
\end{proof}

\subsection{The failure of recentering}

We now prove the main result.

\begin{theorem}[Failure of $\mathsf{FW}(C_\infty,N)$]\label{thm:main}
For every $N\ge 1$, there is no closed family of width~$N$ for
$C_\infty$.
\end{theorem}

\begin{proof}
Suppose for contradiction that $\mathcal{F}$ is a closed family of
width~$N$ for $C_\infty$ with root state $M_\ast$.

\medskip\noindent\textit{Step 1: All reachable states have uniform type
structure.}\;
By Lemma~\ref{lem:type-propagation} and induction on the successor
chain, every state $M$ reachable from $M_\ast$ has
$\exists\Pp.\top\in\tp_M(0_M)$.  The designated $\Pp$-witness for
this formula is some item $w\in I_M$ with $\rho_M(0_M,w)=\Pp$.  By
Lemma~\ref{lem:type-propagation}, $\exists\Pp.\top\in\tp_M(w)$ and
$\forall\Pp.\exists\Pp.\top\in\tp_M(w)$.

Moreover, in the successor state $\Succ(M,w)$: the center has the
same type as $w$ in $M$, which includes $\exists\Pp.\top$ and
$\forall\Pp.\exists\Pp.\top$.  So the analysis repeats.  By
induction, \emph{every state reachable along witness-successor chains
has a center with $\exists\Pp.\top$ in its type}.

\medskip\noindent\textit{Step 2: Every item of a reachable state needs
a $\Pp$-successor.}\;
Actually, we need to be precise.  Not every item needs
$\exists\Pp.\top$ in its type---only items that are $\Pp$ from the
center.  But for the argument it suffices to focus on a single chain.

Start from the root state $M_\ast$.  Its center satisfies
$\exists\Pp.\top$, so it has a $\Pp$-witness $w_1\in I_{M_\ast}$.
Consider $M_1=\Succ(M_\ast,w_1)$.  Its center has type containing
$\exists\Pp.\top$ (by~(R2) and Lemma~\ref{lem:type-propagation}), so
it has a $\Pp$-witness $w_2\in I_{M_1}$.  By~(R1) and~(R3), the
recentering map $u:I_{M_1}\to I_{M_\ast}$ must satisfy:
\begin{align}
u(0_{M_1}) &= w_1, \label{eq:r1}\\
\rho_{M_1}(0_{M_1},w_2) &= \Pp, \label{eq:witness}\\
\rho_{M_\ast}(u(0_{M_1}),u(w_2)) &= \rho_{M_1}(0_{M_1},w_2) = \Pp.
  \label{eq:r3}
\end{align}
Combining \eqref{eq:r1} and \eqref{eq:r3}:
\begin{equation}\label{eq:need}
\rho_{M_\ast}(w_1,u(w_2)) = \Pp.
\end{equation}
So $u(w_2)$ must be an item in $I_{M_\ast}$ that is $\Pp$ from $w_1$.

\medskip\noindent\textit{Step 3: Iterating produces a $\Pp$-chain in
$I_{M_\ast}$.}\;
Set $a_0 = 0_{M_\ast}$ and $a_1 = w_1$.  By~\eqref{eq:need}, $u(w_2)$
is some item $a_2\in I_{M_\ast}$ with $\rho_{M_\ast}(a_1,a_2)=\Pp$.

Now consider $M_2 = \Succ(M_1,w_2) = \Succ(M_\ast,u(w_2)) =
\Succ(M_\ast,a_2)$ (by~(R4)).  The center of $M_2$ has type
containing $\exists\Pp.\top$ (same argument), so it has a
$\Pp$-witness $w_3$.  The recentering
$u':I_{M_2}\to I_{M_\ast}$ satisfies $u'(0_{M_2})=a_2$ and
$\rho_{M_\ast}(a_2,u'(w_3))=\Pp$.  Set $a_3=u'(w_3)$.

Continuing, we obtain a sequence $a_0,a_1,a_2,a_3,\dots$ of items in
$I_{M_\ast}$ with $\rho_{M_\ast}(a_k,a_{k+1})=\Pp$ for all $k\ge 0$.

\medskip\noindent\textit{Step 4: Contradiction.}\;
By Lemma~\ref{lem:spo}, $\Pp$ is a strict partial order on
$I_{M_\ast}$.  In particular, $\Pp$ is antisymmetric and transitive.
For $j<k$ we have $\rho_{M_\ast}(a_j,a_k)=\Pp$ (by transitivity of
the chain), hence $a_j\ne a_k$.  Therefore $a_0,a_1,a_2,\dots$ are
pairwise distinct items in $I_{M_\ast}$.

But $|I_{M_\ast}|\le N$, a finite bound.  An infinite sequence of
pairwise distinct items in a set of size~$N$ is impossible.
Contradiction.
\end{proof}

\begin{remark}
The proof does \emph{not} claim that $C_\infty$ is unsatisfiable.
It~is satisfiable---in infinite models.  The theorem says only that no
\emph{finite closed family of finite-width local states} can capture
$C_\infty$.  Soundness of the tableau (\cite[Corollary~6.6]{GPT2026})
is untouched; it is completeness that fails.
\end{remark}

\begin{remark}
The argument generalizes immediately to any concept $C$ such that
(i)~$C$ is satisfiable only in infinite models, and (ii)~the
infinite model is forced by an unending $\Pp$-chain (or $\Ppi$-chain)
where every element in the chain must have a $\Pp$-successor (or
$\Ppi$-successor) of the same character.  The root cause is the
transitivity of $\Pp$ and $\Ppi$.
\end{remark}

%% ============================================================
\section{Why the obstruction arises and why it is fundamental}
\label{sec:why}

\subsection{The role of strict partial orders}

The recentering axiom~(R3) requires that RCC5 relations in a successor
state are faithfully reflected in the parent state.  When the relation
in question is $\Pp$, this means: if the successor state contains a
$\Pp$-edge, the parent state must contain a corresponding $\Pp$-edge
between the recentered images.

Because $\Pp$ is transitive and antisymmetric, iterating this
requirement along a chain of successor states produces a growing
$\Pp$-chain \emph{inside a single parent state}.  Finiteness of the
state forces the chain to terminate, but the $\exists\Pp$ demands
force it to continue.  This is the contradiction.

\subsection{Comparison with the tree model property}

In standard $\mathcal{ALCI}$ (without RCC constraints), every
satisfiable concept has a tree model.  Tree models can be made
regular (finitely many subtree-types) because there is no interaction
between branches: the tree structure provides natural ``blocking.''

In $\ALCIRCC$, the model is a complete graph.  The $\Pp$-relation
threads through the entire graph and enforces a global partial order.
There is no way to ``block'' a $\Pp$-chain because a blocked node
would need a $\Pp$-successor that is \emph{strictly deeper} in the
partial order---reusing an ancestor is impossible (it would reverse
$\Pp$).  This is the fundamental obstacle that distinguishes
$\ALCIRCC$ from standard description logics.

\subsection{Comparison with the quasimodel approach}

The quasimodel method of \cite{Claude2026} has a complementary gap.
There, \emph{soundness} (satisfiable concept $\Rightarrow$
quasimodel) is proven, but \emph{completeness} (quasimodel
$\Rightarrow$ model) has a gap: the Henkin construction may fail for
abstract quasimodels.

The contextual tableau of \cite{GPT2026} has the opposite gap:
\emph{soundness} (closed family $\Rightarrow$ model) is proven, but
\emph{completeness} (satisfiable concept $\Rightarrow$ closed family)
fails outright for concepts like $C_\infty$.

\begin{center}
\begin{tabular}{lcc}
\toprule
& Quasimodel~\cite{Claude2026} & Contextual tableau~\cite{GPT2026} \\
\midrule
Satisfiable $\Rightarrow$ representation & Proven & \textbf{False} for $C_\infty$ \\
Representation $\Rightarrow$ model & \textbf{Gap} & Proven \\
\bottomrule
\end{tabular}
\end{center}

Neither approach, as formulated, yields a decision procedure.

%% ============================================================
\section{Possible directions forward}\label{sec:directions}

The failure of $\mathsf{FW}(C,N)$ does not settle the decidability
question.  It eliminates one specific proof strategy (fixed-width
closed families) but leaves several alternatives open.

\begin{enumerate}
\item \textbf{Variable-width or infinite-width representations.}\;
    The finite-width constraint is the source of the obstruction.
    One could allow the width of local states to depend on the
    \emph{unfolding depth} rather than being globally bounded.  The
    challenge is ensuring that the resulting search space remains
    finite (or decidable).

\item \textbf{Automata-theoretic approaches.}\;
    Two-way alternating tree automata handle $\mathcal{ALCI}$
    decidability elegantly.  Adapting them to the complete-graph
    setting of $\ALCIRCC$ is non-trivial, but the regular structure
    of the infinite $\Pp$-chain (every node has the same type)
    suggests that some form of $\omega$-regularity or
    B\"uchi acceptance might capture the required infinite patterns.

\item \textbf{Reduction to constraint satisfaction.}\;
    The quasimodel approach of~\cite{Claude2026} reduces the problem
    to checking whether abstract quasimodels (which can be finitely
    enumerated) yield genuine models.  Closing the extension gap in
    that approach---showing that every abstract quasimodel
    satisfying~(Q1)--(Q3) is realizable---would still yield decidability.
    The present counterexample does not affect that direction.

\item \textbf{Undecidability.}\;
    It is also possible that $\ALCIRCC$ is undecidable.  The
    $\Pp$-chain phenomenon and the absence of finite model or tree
    model properties are suggestive, though not conclusive.  A
    reduction from a known undecidable problem (e.g., the halting
    problem or the tiling problem) to $\ALCIRCC$ satisfiability would
    settle the question negatively.
\end{enumerate}

%% ============================================================
\section{The counterexample for RCC8}\label{sec:rcc8}

The same argument applies to $\mathsf{ALCI\_RCC8}$.  In RCC8, the
role $\mathsf{NTPP}$ (non-tangential proper part) is a strict partial
order, with composition
$\mathsf{NTPP}\circ\mathsf{NTPP}=\{\mathsf{NTPP}\}$.  The concept
\[
C'_\infty = (\exists\mathsf{NTPP}.\top)
  \sqcap(\forall\mathsf{NTPP}.\exists\mathsf{NTPP}.\top)
\]
forces an infinite $\mathsf{NTPP}$-chain, and the strict partial order
argument of Theorem~\ref{thm:main} applies verbatim.  So
$\mathsf{FW}(C'_\infty,N)$ fails for every~$N$ in the RCC8 setting as
well.

%% ============================================================
\section{Conclusion}

We have shown that the finite-width extraction conjecture
$\mathsf{FW}(C,N)$ is false for the satisfiable $\ALCIRCC$ concept
$C_\infty=(\exists\Pp.\top)\sqcap(\forall\Pp.\exists\Pp.\top)$.  The
proof is elementary: transitivity of $\Pp$ and the recentering
axiom~(R3) together force an infinite $\Pp$-chain inside any
candidate local state, contradicting the finite width bound.

This closes the completeness direction of the contextual tableau
approach of~\cite{GPT2026} negatively.  Combined with the extension
gap in the quasimodel approach of~\cite{Claude2026}, the decidability
of $\ALCIRCC$ (and $\mathsf{ALCI\_RCC8}$) remains an open problem.
Both approaches have been shown to have genuine, non-trivial gaps that
are not mere artifacts of incomplete proofs but reflect fundamental
structural features of the logic.

The key structural feature is the combination of (i)~the transitivity
of $\Pp$ (creating strict partial orders), (ii)~the universal
propagation of $\forall\Pp$ along $\Pp$-chains, and (iii)~the
requirement that models be complete graphs.  Any successful
decidability proof---or undecidability reduction---for $\ALCIRCC$ will
need to engage with this combination directly.

\begin{thebibliography}{9}

\bibitem{Wessel2003}
Michael Wessel.
\newblock \emph{Qualitative Spatial Reasoning with the ALCIRCC Family
  -- First Results and Unanswered Questions}.
\newblock Technical Report FBI-HH-M-324/03, University of Hamburg,
  2002/2003.

\bibitem{Claude2026}
Claude.
\newblock \emph{On the Decidability of ALCIRCC5 and ALCIRCC8:
  Quasimodels Meet the Patchwork Property}.
\newblock Discussion draft, 2026.

\bibitem{GPT2026}
GPT-5.4 Pro.
\newblock \emph{A Contextual Tableau Calculus for $\ALCIRCC$}.
\newblock Discussion draft, 2026.

\bibitem{RenzNebel1999}
Jochen Renz and Bernhard Nebel.
\newblock On the complexity of qualitative spatial reasoning:
  A~maximal tractable fragment of the Region Connection Calculus.
\newblock \emph{Artificial Intelligence}, 108(1--2):69--123, 1999.

\bibitem{Renz1999}
Jochen Renz.
\newblock Maximal tractable fragments of the Region Connection
  Calculus: A~complete analysis.
\newblock In \emph{Proc.\ IJCAI}, pages 448--455, 1999.

\end{thebibliography}

\end{document}
